\documentclass[11pt,a4paper]{article}

% ---------- packages ----------
\usepackage[utf8]{inputenc}
\usepackage[margin=2.5cm]{geometry}
\usepackage{microtype}
\usepackage{array}
\usepackage{booktabs}
\usepackage{tabularx}
\usepackage{longtable}
\usepackage{enumitem}
\usepackage[table]{xcolor}
\usepackage{caption}
\usepackage[numbers,sort&compress]{natbib}
\usepackage{xurl}
\usepackage[colorlinks=true,linkcolor=blue!50!black,citecolor=blue!50!black,urlcolor=blue!50!black]{hyperref}

\setlist{itemsep=2pt,topsep=4pt}
\setlength{\parskip}{4pt}
\renewcommand{\arraystretch}{1.25}

% one source = one subsection with a short label/value list
\newenvironment{sourceinfo}{\begin{description}[leftmargin=0pt,labelsep=0.5em,font=\normalfont\bfseries]}{\end{description}}
\newcommand{\source}[2]{\subsection{\texorpdfstring{#1~\cite{#2}}{#1}}}

% ---------- title ----------
\title{\textbf{Literature Review}\\[6pt]
\large Development of an Intelligent Kanban System with an AI Assistant\\
for Automatic Task Generation from Voice Messages}
\author{Symbat Saparbay | Akbope Bakytkeldy\\ \small YouTrack: DEMO-20}
\date{}

\begin{document}
\maketitle

% =====================================================
\section{Introduction}

Existing task management systems usually require users to type out task titles, descriptions, dates and other attributes manually. The proposed system is based on a different approach: the user records a voice message, an automatic speech recognition (ASR) model produces a transcript, an AI component finds the actionable information in it, and the backend validates a draft task that the user can review before it is added to the Kanban board.

The literature is structured according to the three search techniques the team used to collect the papers: (1)~broad concept-based searches, (2)~narrow context-, method- or language-based searches, and (3)~searches for adjacent techniques and technical issues. These labels describe how \emph{the team conducted its search}; they are not classifications made by the authors of the papers.

% =====================================================
\section{Research Questions}

\begin{enumerate}[label=\textbf{RQ\arabic*.}, leftmargin=*]
  \item How can AI support task management and workflow automation?
  \item Which ASR approaches merit evaluation for voice messages, especially in Kazakh?
  \item How can task intent and fields such as title, due date and assignee be extracted from spoken language?
  \item How can a backend avoid storing incomplete or fabricated task information?
  \item How should accuracy and reliability be evaluated before a Kanban card is created?
\end{enumerate}

% =====================================================
\section{Search Strategy and Scope}

The team used three kinds of queries:
\begin{itemize}
  \item \textbf{Type 1 -- Broad terms:} Kanban in software and project work, AI-assisted task management, NLP-based project management.
  \item \textbf{Type 2 -- Narrow terms or methods:} self-supervised and large-scale ASR, Kazakh speech corpora and ASR, spoken task-oriented dialogue, intent detection, slot filling, action-item extraction in meetings, structured LLM output.
  \item \textbf{Type 3 -- Related terms and contradictions:} API argument filling, structured hallucinations.
\end{itemize}

\textbf{Filtering strategy.} Candidate publications were filtered by research problem, proposed method, reported results, limitations and possible application in our project. This is a narrative, not a systematic, literature review. Fourteen sources were retained.

\textbf{Source quality and recency.} Bibliographic metadata of every retained source was checked against the publisher page, DOI record or proceedings on 23 October 2026. Three sources from the first draft were removed because their venue or publication record could not be confirmed (a multi-agent workflow paper, a review of action-item extraction and a general ASR/NLP overview); the action-item review was replaced with a peer-reviewed ACM conference paper on the same problem. Twelve of the fourteen sources are peer-reviewed journal articles or main-conference papers; one is an ICML workshop paper and one is a journal article from a smaller, non-indexed venue, and both are used only as supporting evidence. Twelve sources were published in 2021--2026. Two foundational works are retained deliberately: the Kanban mapping study~\cite{ahmad2018kanban} and wav2vec~2.0~\cite{baevski2020wav2vec}, because the board model and the self-supervised ASR models evaluated for Kazakh build on them.

% =====================================================
\section{Type 1 -- Broad Searches: AI and Task-Management Systems}

\source{Ahmad et al. (2018): Kanban in Software Engineering: A Systematic Mapping Study}{ahmad2018kanban}
\begin{sourceinfo}
  \item[Problem.] Kanban is widely adopted by software teams, but the evidence on its benefits and difficulties is scattered.
  \item[Method.] A systematic mapping study of published research on Kanban in software engineering, classifying studies by topic, research type and reported outcomes.
  \item[Reported findings.] The mapped studies associate Kanban with visualizing work, limiting work in progress and improving coordination, but also report adoption difficulties and a limited amount of rigorous empirical evidence.
  \item[Relevance.] Grounds the board model of our system (columns, visible status, WIP) in established practice rather than in a single product, and supports keeping the board itself conventional while the novelty lies in task entry.
  \item[Limitation for our project.] Foundational and published in 2018; it concerns software teams and does not address AI or voice input, so it motivates the board design but not the assistant.
\end{sourceinfo}

\source{Adebanjo et al. (2026): Design and Implementation of an AI-Powered Task Management System}{adebanjo2026taskwise}
\begin{sourceinfo}
  \item[Problem.] Manual task entry and the lack of intelligent assistance make task management less convenient.
  \item[Method.] TaskWise combines natural language processing, scheduling and reminders, a React/Tailwind front end, and Firebase/Firestore back-end services with synchronization.
  \item[Findings.] The paper reports usability improvements from task automation and natural language input.
  \item[Relevance.] A close analogue of our product; useful when discussing user-friendly task creation and quick feedback after natural language input.
  \item[Limitations.] Dependence on online Firebase services; rule- and context-based AI rather than learning from history; no collaborative workspace; web version only; small student test group. Published in a smaller, non-indexed journal, so it is used as a product analogue rather than as empirical evidence.
  \item[Our difference.] We focus specifically on the \emph{speech $\rightarrow$ transcription $\rightarrow$ task extraction $\rightarrow$ user confirmation $\rightarrow$ Kanban card} pipeline and evaluate each step separately.
\end{sourceinfo}

\source{Chen et al. (2021): Smart Project Management: Interactive Platform Using NLP Technology}{chen2021smart}
\begin{sourceinfo}
  \item[Problem.] Communication and document-processing difficulties among project stakeholders.
  \item[Method.] Construction-specific contract keyword detection using a pipeline of PSENet, CRNN and BRNN-CNN architectures.
  \item[Reported findings.] The authors present a platform supporting real-time document processing and communication in construction contract management.
  \item[Relevance.] Shows how domain-specific NLP can reduce manual extraction from unstructured project data.
  \item[Limitation for our project.] Based on text and document analysis of construction contracts rather than speech commands and Kanban task generation; the architecture should not be adopted directly.
\end{sourceinfo}

% =====================================================
\section{Type 2 -- Narrow Searches: Speech Understanding and Structured Task Extraction}

\source{Weld et al. (2022): A Survey of Joint Intent Detection and Slot Filling Models}{weld2022survey}
\begin{sourceinfo}
  \item[Problem.] A user request contains both an overall intent and the values needed to carry it out.
  \item[Approach.] A survey of joint intent classification and slot filling: metrics, datasets and deep learning architectures.
  \item[Findings.] Emphasizes the strong link between intent recognition and semantic slot extraction.
  \item[Relevance.] Provides the theoretical basis for detecting a \texttt{create\_task} intent and extracting slots such as \texttt{title}, \texttt{due\_date} and \texttt{assignee}.
  \item[Limitation.] A survey of existing models and experiments, not an implementation or benchmark for a Kanban setting; some methods are dated compared with modern LLMs.
\end{sourceinfo}

\source{Baevski et al. (2020): wav2vec 2.0: Self-Supervised Learning of Speech Representations}{baevski2020wav2vec}
\begin{sourceinfo}
  \item[Problem.] Accurate ASR normally requires large amounts of transcribed speech, which low-resource languages such as Kazakh lack.
  \item[Method.] Self-supervised pre-training on unlabelled audio by masking latent speech representations and solving a contrastive task, followed by fine-tuning on a small labelled set.
  \item[Reported findings.] With ten minutes of labelled data and pre-training on 53k hours of unlabelled audio, the model reaches 4.8/8.2 WER on the LibriSpeech clean/other test sets.
  \item[Relevance.] Foundational work behind the wav2vec2 and XLS-R models evaluated for Kazakh~\cite{kozhirbayev2023kazakh}; explains why a Kazakh-specific fine-tuned model is a realistic candidate for our STT benchmark.
  \item[Limitation.] Results are for read English speech; fine-tuning and evaluation on Kazakh and mixed speech are still required.
\end{sourceinfo}

\source{Radford et al. (2023): Robust Speech Recognition via Large-Scale Weak Supervision (Whisper)}{radford2023whisper}
\begin{sourceinfo}
  \item[Problem.] ASR models fine-tuned on one dataset often degrade on other recording conditions and domains.
  \item[Method.] Whisper: an encoder--decoder Transformer trained on 680{,}000 hours of weakly supervised multilingual and multitask audio, including language identification and translation.
  \item[Reported findings.] Without dataset-specific fine-tuning, the models are often competitive with fully supervised systems on standard benchmarks and approach human robustness.
  \item[Relevance.] The base multilingual model for Russian and English and the automatic language-detection step in our STT routing (BE-03).
  \item[Limitation.] Accuracy varies strongly by language and depends on the amount of training data per language; zero-shot performance on Kazakh and on code-switched Kazakh--Russian speech must be measured rather than assumed.
\end{sourceinfo}

\source{Mussakhojayeva, Khassanov and Varol (2022): KSC2: An Industrial-Scale Open-Source Kazakh Speech Corpus}{mussakhojayeva2022ksc2}
\begin{sourceinfo}
  \item[Problem.] Progress in Kazakh ASR is limited by the lack of large, openly available transcribed speech.
  \item[Method.] KSC2: more than a thousand hours of manually transcribed Kazakh speech, about three times the size of the original KSC, validated with preliminary ASR experiments on several evaluation sets.
  \item[Reported findings.] The corpus is released openly and supports training and evaluation of Kazakh ASR models.
  \item[Relevance.] A source of training data and reference test sets for Kazakh models in our STT benchmark, alongside our own recorded commands.
  \item[Limitation.] Mostly monolingual Kazakh; it does not represent short workplace commands or Kazakh--Russian code-switching typical of our users.
\end{sourceinfo}

\source{Kozhirbayev (2023): Kazakh Speech Recognition: Wav2vec2.0 vs.\ Whisper}{kozhirbayev2023kazakh}
\begin{sourceinfo}
  \item[Problem.] Kazakh speech recognition is difficult because annotated data are limited.
  \item[Method.] Comparison of wav2vec2.0, wav2vec2-XLS-R and Whisper for Kazakh speech recognition.
  \item[Reported findings.] wav2vec2.0 achieves a WER of 8.7 and a CER of 2.8 on the authors' test set, and large Whisper reaches a CER of approximately 4. Metrics and configurations should be checked before generalizing the model ranking.
  \item[Relevance.] A priority paper for selecting and testing ASR models for Kazakh voice messages.
  \item[Limitation.] Results on the paper's dataset do not guarantee the same performance with student microphones, informal speech, noisy conditions or mixed Kazakh/Russian speech.
\end{sourceinfo}

\source{Si et al. (2023): SpokenWOZ: A Speech-Text Benchmark for Spoken Task-Oriented Dialogue}{si2023spokenwoz}
\begin{sourceinfo}
  \item[Problem.] Written dialogue benchmarks may miss the challenges of real spoken language.
  \item[Method.] The SpokenWOZ dataset and experiments on spoken task-oriented dialogue: about 5{,}700 dialogues, 203{,}000 turns and 249 hours of audio across eight domains.
  \item[Reported findings.] Current systems struggle with spoken dialogue, including dialogue state tracking and task completion.
  \item[Relevance.] Motivates evaluation on real voice samples with disfluencies, transcription errors and multi-step instructions rather than typed text only.
  \item[Limitation.] Targets conversational dialogue, not card generation, and does not include Kazakh.
\end{sourceinfo}

\source{Golia and Kalita (2023): Action-Item-Driven Summarization of Long Meeting Transcripts}{golia2023action}
\begin{sourceinfo}
  \item[Problem.] Meeting transcripts mix conversation, decisions and commitments to action, and long transcripts exceed what a language model can track reliably.
  \item[Method.] The transcript is split into topic-based sections; action items are extracted from each section, sections are summarized recursively and in parallel, and the results are merged into an action-item-driven summary. Three section-splitting methods are compared.
  \item[Reported findings.] On the AMI meeting corpus the pipeline reaches a BERTScore of 64.98, about 4.98\% above a fine-tuned BART baseline.
  \item[Relevance.] Supports extracting several separate action items from one input and processing long meeting input in sections, which corresponds to our multiple-task and meeting-summary requirements.
  \item[Limitation.] English meeting transcripts and a summary-quality metric; it does not measure whether each extracted action item has the correct assignee and deadline, which our evaluation must check field by field.
\end{sourceinfo}

\source{Sun et al. (2024): Speech-based Slot Filling using Large Language Models}{sun2024speech}
\begin{sourceinfo}
  \item[Problem.] ASR errors can degrade downstream value extraction from voice commands.
  \item[Method.] LLM-based speech slot filling using prompting, LoRA-based noise-robust fine-tuning and linearized knowledge injection, evaluated on SLURP.
  \item[Reported findings.] The proposed Vicuna-based model improves SLU-F1 over Flan-T5-XL baselines in limited-data and zero-shot settings.
  \item[Relevance.] Highly relevant: ASR output contains errors, so value extraction must be robust to them.
  \item[Limitation.] The dataset and models differ from our Kanban use case.
\end{sourceinfo}

\source{Li, Ramprasad and Zhang (2024): Improving Structured LLM Output for Information Extraction}{li2024simple}
\begin{sourceinfo}
  \item[Problem.] LLM outputs are unreliable when an exact structured format is required.
  \item[Method.] Generate-and-Organize (G\&O): the model first generates content and then organizes it into the target structure; evaluated on named entity recognition and relation extraction.
  \item[Reported findings.] The authors report improved structured information extraction.
  \item[Relevance.] One way to obtain task JSON; regardless of the prompting strategy, the backend must still perform schema validation.
  \item[Limitation.] NER and relation extraction do not correspond exactly to identifying tasks and relative deadlines in voice transcripts.
\end{sourceinfo}

% =====================================================
\section{Type 3 -- Adjacent Solutions and Contradictions}

\source{Mok et al. (2024): LLM-based Frameworks for API Argument Filling}{mok2024llm}
\begin{sourceinfo}
  \item[Problem.] Conversational LLM output may not conform to the required schema or to the user's exact wording.
  \item[Solution.] Grounding together with prompting and training techniques for API argument filling.
  \item[Reported results.] Positive results for grounding methods in API argument filling.
  \item[Relevance.] An extracted task must match a predefined backend schema and faithfully reflect the user's request, which matters for correct dates, names and task attributes.
  \item[Limitation.] Addresses task-oriented API argument filling, not Kanban services or our target languages.
\end{sourceinfo}

\source{Hoang (2026): Structured Hallucination in Tool-Using Agents}{hoang2026structured}
\begin{sourceinfo}
  \item[Problem.] An LLM can silently corrupt structured values even when its response is syntactically valid JSON or Markdown.
  \item[Method.] A production-derived benchmark of 200 runs of LLM synthesis of travel API results, covering different formats and the use of SSE metadata as a transport mechanism.
  \item[Reported findings.] In this test, LLM-mediated output formatting corrupted structured fields; avoiding unnecessary re-synthesis preserved tool values.
  \item[Relevance.] A useful caution that a valid JSON object does \emph{not} guarantee correct field values. The transcript should be preserved, the output validated, and uncertain fields confirmed by the user.
  \item[Limitation.] A single deployment and model, focused on reformatting structured tool outputs rather than extracting information from audio; the error rates do not generalize to our use case.
\end{sourceinfo}

% =====================================================
\section{Comparison of Sources}

\begingroup
\small
\setlength{\tabcolsep}{4pt}
\rowcolors{2}{gray!8}{white}
\begin{longtable}{>{\raggedright\arraybackslash}p{0.6cm} >{\raggedright\arraybackslash}p{1.5cm} >{\raggedright\arraybackslash}p{2.5cm} >{\raggedright\arraybackslash}p{3.0cm} >{\raggedright\arraybackslash}p{3.1cm} >{\raggedright\arraybackslash}p{3.0cm}}
\caption{Comparison of the fourteen reviewed sources.}\label{tab:comparison}\\
\toprule
\textbf{No.} & \textbf{Search type} & \textbf{Publication} & \textbf{Main approach} & \textbf{Value for our project} & \textbf{Applicability limit} \\
\midrule
\endfirsthead
\toprule
\textbf{No.} & \textbf{Search type} & \textbf{Publication} & \textbf{Main approach} & \textbf{Value for our project} & \textbf{Applicability limit} \\
\midrule
\endhead
\cite{ahmad2018kanban} & Broad & Ahmad et al. (2018) & Systematic mapping of Kanban research & Board model, WIP and visible status & Foundational; no AI or voice \\
\cite{adebanjo2026taskwise} & Broad & Adebanjo et al. (2026) & NLP-enabled TaskWise platform & Comparable task-management product & Small evaluation; smaller venue \\
\cite{chen2021smart} & Broad & Chen et al. (2021) & NLP and contract keyword detection & Domain-focused information extraction & Construction documents \\
\cite{weld2022survey} & Narrow & Weld et al. (2022) & Survey of intent and slot models & Intent and task-field design & Not a task-generation experiment \\
\cite{baevski2020wav2vec} & Narrow & Baevski et al. (2020) & Self-supervised ASR pre-training & Basis of Kazakh wav2vec2/XLS-R models & Foundational; English read speech \\
\cite{radford2023whisper} & Narrow & Radford et al. (2023) & Whisper, weakly supervised multilingual ASR & Base model and language detection & Per-language accuracy varies \\
\cite{mussakhojayeva2022ksc2} & Narrow & Mussakhojayeva et al. (2022) & KSC2 Kazakh speech corpus & Kazakh training and test data & No code-switching or commands \\
\cite{kozhirbayev2023kazakh} & Narrow & Kozhirbayev (2023) & Kazakh ASR model comparison & Candidate ASR models and WER/CER & Domain and recording conditions vary \\
\cite{si2023spokenwoz} & Narrow & Si et al. (2023) & SpokenWOZ benchmark & Realistic speech evaluation cases & Different domains and languages \\
\cite{golia2023action} & Narrow & Golia \& Kalita (2023) & Action-item-driven meeting summarization & Several action items per input & English; summary metric, not fields \\
\cite{sun2024speech} & Narrow & Sun et al. (2024) & LLM slot filling under noisy ASR & Robust extraction from transcripts & Different dataset and models \\
\cite{li2024simple} & Narrow & Li et al. (2024) & Generate-and-Organize (G\&O) & Structured output strategy & Tested on NER/RE \\
\cite{mok2024llm} & Adjacent & Mok et al. (2024) & Grounded API argument filling & Faithfulness to schema and user text & Not Kanban-specific \\
\cite{hoang2026structured} & Adjacent & Hoang (2026) & Structured hallucination benchmark & Field fidelity and validation risks & Travel-data synthesis, not STT \\
\bottomrule
\end{longtable}
\endgroup

% =====================================================
\section{Synthesis: Implications for the Proposed Solution}

\paragraph{A modular pipeline is appropriate.}
The reviewed sources treat speech recognition, spoken-language understanding, task and action-item extraction, and reliable structured output as separable problems. Our proposed architecture is therefore:
\begin{center}
\small
Voice message $\rightarrow$ ASR transcription $\rightarrow$ task/intent extraction $\rightarrow$ schema validation $\rightarrow$ user review $\rightarrow$ backend storage $\rightarrow$ Kanban card
\end{center}

\paragraph{Model selection is empirical.}
Whisper~\cite{radford2023whisper} provides a robust multilingual baseline with language detection, while self-supervised models~\cite{baevski2020wav2vec} fine-tuned on Kazakh data such as KSC2~\cite{mussakhojayeva2022ksc2} are the main alternative for Kazakh. The Kazakh ASR comparison~\cite{kozhirbayev2023kazakh} motivates testing both families on our own recordings. Universal claims of model superiority should not be made on the basis of a single dataset, and none of these sources covers Kazakh--Russian code-switching.

\paragraph{Task extraction must be evaluated independently of transcription.}
The intent detection, slot filling and spoken-dialogue papers~\cite{weld2022survey,si2023spokenwoz,sun2024speech} show that even a good transcript can lead to wrongly extracted entities or actions. Both stages need to be assessed.

\paragraph{Valid JSON syntax is not enough.}
The structured extraction, API argument filling and hallucination papers~\cite{li2024simple,mok2024llm,hoang2026structured} imply two complementary checks: (a)~\emph{schema validity} (types, required fields, valid statuses) and (b)~\emph{semantic correctness} with respect to the transcript. Schema validation must be performed explicitly before data are stored.

\paragraph{The MVP should stay small.}
The Kanban literature~\cite{ahmad2018kanban} supports keeping the board itself conventional, and comparable task-management products~\cite{adebanjo2026taskwise} show that the difficult part is reliable task entry rather than the board. A basic backend for ASR, task extraction and validation is easier to justify and test than autonomous multi-agent workflow optimization, which is outside the MVP.

% =====================================================
\section{Proposed Evaluation Plan}

\begingroup
\small
\rowcolors{2}{gray!8}{white}
\begin{tabularx}{\textwidth}{>{\raggedright\arraybackslash}p{2.6cm} >{\raggedright\arraybackslash}X >{\raggedright\arraybackslash}X}
\toprule
\textbf{Component} & \textbf{Proposed measurement} & \textbf{What to check} \\
\midrule
ASR & Word Error Rate (WER); optionally CER & Kazakh and Russian recordings; background noise \\
Intent & Precision, recall, F1 & Task requests correctly distinguished from ordinary speech \\
Task extraction & Field-level precision/recall/F1 or exact match & Title, deadline, assignee, description \\
Structured output & Schema-valid response rate & Correct data types, required fields, date formats \\
Faithfulness & Unsupported-value / fabricated-field rate & Did the AI invent a deadline the user never mentioned? \\
End-to-end & Correct task creation rate; user correction rate; latency & Can the user confirm and save the intended Kanban card? \\
\bottomrule
\end{tabularx}
\endgroup

\medskip
Evaluation should use a small labelled set of representative voice messages: single-task commands, several tasks in one message, unclear dates, missing assignees, Kazakh and Russian messages, and messages that contain no actionable task. Final metrics and thresholds will be agreed with the team and the supervisor.

% =====================================================
\section{Conclusion}

The fourteen reviewed publications cover complementary parts of the intelligent Kanban problem. Kanban and task-management studies establish the relevance of the product; ASR, Kazakh speech and spoken-dialogue studies highlight language and audio challenges; intent detection, slot filling and structured extraction studies inform the AI pipeline; and API-grounding and hallucination studies identify reliability risks. Together they \emph{motivate}, but do not directly validate, the planned voice-to-Kanban workflow. The next step is to select candidate tools, build a baseline implementation and measure its performance on representative user utterances before finalizing the technology stack.

% =====================================================
\bibliographystyle{unsrtnat}
\bibliography{BE01}

\end{document}